% =====================================================================
%  DRAFT -- Re-embed, Map, or Window?
%  The Economics of Embedding Model Upgrades
%
%  Rewritten 2026-11-05 for framing B (RUNBOOK Phase 5.9, DECISIONS.md).
%  The A-framing original is kept beside this file as
%  reuse_strategies_draft_A_original.tex for diffing.
%
%  NUMBERS: every figure in this draft is tier PAPER in ../PAPER_NUMBERS.md,
%  recomputed from the result JSONs of the 2026-11-04 1M run. \todo{} now marks
%  only what is genuinely unmeasured -- chiefly the K' = 10/25 frontier
%  (RUNBOOK Phase 11.1), which is the one open item that moves a headline number.
%  Target: see paper/OUTLINE.md. ICDE 2027 offers Research, DEFT and PhD Symposium only --
%  there is no experiments-and-analysis track there, so the framing cannot lean on one.
% =====================================================================
\documentclass[conference]{IEEEtran}
\IEEEoverridecommandlockouts
\usepackage{amsmath,amssymb,amsfonts}
\usepackage{algorithmic}
\usepackage{graphicx}
\usepackage{booktabs}
\usepackage{multirow}
\usepackage{xcolor}
\usepackage{url}
\usepackage[hidelinks]{hyperref}

\newcommand{\todo}[1]{\textcolor{red}{[TODO: #1]}}
\newcommand{\note}[1]{\textcolor{blue}{[NOTE: #1]}}
\newcommand{\topk}{\mathrm{Top}}
\newcommand{\win}{W}
\newcommand{\Gk}{G(K')}
\newcommand{\cov}{\mathrm{cov}}
\newcommand{\debt}{\mathrm{debt}}

\begin{document}

\title{Re-embed, Map, or Window?\\The Economics of Embedding Model Upgrades}

\author{\IEEEauthorblockN{Anonymous Authors}
\IEEEauthorblockA{\textit{Under double-blind review}}}

\maketitle

% ---------------------------------------------------------------------
\begin{abstract}
Upgrading the embedding model behind a vector index normally means re-encoding the
whole corpus before the new model can serve. Several reuse strategies avoid or defer
that cost: mapping new-model queries into the old space with a learned adapter,
re-ranking a window of candidates nominated by the old index with the new model, and
combinations of these with lexical nomination and background migration. Each has been
evaluated in isolation, on its own metric, and mostly on text. We do not propose
another strategy. We evaluate all of them under one cost model across twelve model
pairs on a 1M-passage text corpus and a 45K-document code corpus, charging each
strategy for the work it has done at the moment it first serves each query.

We set out to build a \emph{regime map}: two cheap pair-level measurements --- space
compatibility $\kappa$ and window residual miss $\rho(K')$ --- predicting which
strategy to use. That hypothesis was pre-registered with a falsification condition, and
it failed. No selector built on those axes beats a constant predictor; the constant's
mean regret, $0.0019$ Recall@10, falls inside the equivalence band we fixed before
running anything; and at the traffic intensity we pre-registered, the background migration
order leaves served quality unchanged. Three of our five pre-registered hypotheses were
falsified on their own stated criteria, and we report each as it stands.

What the measurements do support is an economic account. Mapping never wins: across
336 pair\,$\times$\,setting\,$\times$\,tolerance configurations no adapter-bridge
strategy is ever selected, and bridging costs between $0.12$ and $0.82$ Recall@10
against the target model. And the cost of windowed reuse can be bounded in advance without
simulation. A closed form assuming independent windows, $1 - (1 - K'/N)^{|Q|}$, overstates
that cost by $10$--$24\%$ across the operating range; the union of the windows a workload
actually touches, obtained by retrieval alone, is tighter, and the gap between them is exactly
the window concentration. The practical upgrade decision is therefore not \emph{which}
reuse strategy, but \emph{whether} the corpus and workload sit in the regime where
windowed reuse pays at all --- a question answerable in advance, in closed form, from
one static measurement.
All code, configurations, and per-pair measurements are released.
\end{abstract}

\begin{IEEEkeywords}
vector search, embedding models, index maintenance, model migration, benchmark,
negative results
\end{IEEEkeywords}

% =====================================================================
\section{Introduction}
% =====================================================================

Dense retrieval systems tie their index to a specific embedding model. When a better
model appears, the vectors in the index become incompatible with the new query
encoder, and the standard practice is a blue--green rebuild: re-encode the corpus,
build a second index, validate it, and switch traffic~\cite{zilliz-migration}. At the
scale of hundreds of millions of documents this costs days to weeks of accelerator
time before the new model serves a single query.

Three families of reuse strategies have emerged to avoid that up-front cost.
\emph{Space mapping} learns a light transformation that maps new-model queries into the
old space so the existing index keeps serving~\cite{drift-adapter}.
\emph{Window re-ranking} retrieves a bounded candidate window from the old index,
scores it with the new model, and materializes new vectors progressively in the
background~\cite{embedflow,thaw-anon}. \emph{Structure reuse} exploits the observation
that old-model neighborhoods remain partially informative under the new model, for
instance to accelerate building the new graph index~\cite{dgm}.

These strategies were evaluated on different metrics, different model pairs, and
almost exclusively on text. A practitioner facing a concrete upgrade has no basis for
choosing among them, and the published numbers cannot be compared. Moreover, space
mapping does not upgrade stored vectors: it bounds retrieval by what the old model
encoded, so it is a bridge, not a destination. Any fair comparison must measure the
cost of \emph{reaching the new model} while keeping quality acceptable along the way.

\textbf{What we expected, and what we found.} We designed this study around a regime
map: if $\kappa$ and $\rho(K')$ predict the winning strategy on a held-out pair, a
practitioner can measure two cheap quantities and read off an answer. We pre-registered
that hypothesis with an explicit falsification condition --- the $(\kappa,\rho)$ rule's
leave-one-pair-out mean regret must beat a majority-class baseline's --- and ran it.
The baseline won. Two further pre-registered hypotheses also failed on their stated
criteria. Rather than retune, we report the design as it was registered and the
outcomes as they came out; Section~\ref{sec:hyp} is a table of all five with their
verdicts.

The same measurements, however, settle the practitioner's question by a different
route. The reason no selector helps is that there is almost nothing to select: one
strategy family dominates everywhere we measured, the variants within it are
statistically indistinguishable, and the quantity that actually varies --- what the
migration costs --- can be bounded from the source index alone, with no replay and no target
encoding.

\textbf{Findings.}
\begin{itemize}
\item \textbf{F1 (mapping does not work).} Across 336
pair\,$\times$\,setting\,$\times$\,tolerance configurations, no adapter-bridge strategy
(S1) and no mapped-window strategy (S3) is ever selected. Bridged retrieval serves
between $0.05$ and $0.76$ Recall@10 where the target model serves $0.83$--$0.88$, a
shortfall of $0.12$ to $0.82$ (Section~\ref{sec:cost}).
\item \textbf{F2 (window cost admits a tight static bound).} Encoding the union of the
served windows is a lower bound on what windowed reuse must pay, by construction. The
question is how tight it is, and against what. A closed form needing no retrieval
($1 - (1 - K'/N)^{|Q|}$) overstates cost by $10.6$--$15.5\%$ at $K' = 50$ and
$16.6$--$24.3\%$ at $K' = 200$, because real windows overlap more than independent draws;
the measured union closes that gap, and the replay confirms the bound is attained rather
than exceeded by background work (Section~\ref{sec:econ}). This is the result we would hand
a practitioner --- as a provisioning bound, not a prediction.
\item \textbf{F3 (selection is moot).} No feature-based rule beats a constant
predictor; the two rules that match its accuracy reproduce its regret to every digit;
and the constant's mean regret falls inside the pre-registered equivalence band
(Section~\ref{sec:select}).
\item \textbf{F4 (the regime boundary is the real decision).} Windowed reuse needs
$|Q|\,K'$ small against $N$. On our code corpus, whose query set approaches its
document count, window coverage is $1.0000$ at every depth and the ceiling on savings
is exactly zero (Section~\ref{sec:regime}).
\end{itemize}

\textbf{Contributions.} (i) A unified evaluation frame that charges every reuse
strategy for the encoding work it has done when it first serves each query, with
completion of the migration as an admissibility constraint rather than a cost
(Section~\ref{sec:frame}). (ii) A closed-form, replay-free cost model for windowed
reuse, validated to four decimal places (Section~\ref{sec:econ}). (iii) A measured
comparison of all three strategy families over twelve pairs on text and code
(Section~\ref{sec:cost}). (iv) Three pre-registered negative results, reported with
their criteria and their sensitivity analyses (Sections~\ref{sec:select}--\ref{sec:order}).

\textbf{Scope.} We study the \emph{complete-migration} regime: every strategy we admit
ends with the corpus encoded under the target model. We make no exactness guarantees
for mixed-state retrieval, propose no new serving mechanism, and do not study the
economics of deliberately incomplete migration; those are treated in concurrent
work~\cite{thaw-anon}. Section~\ref{sec:disc} states what this scope costs us.

% =====================================================================
\section{Problem Setting and Evaluation Frame}\label{sec:frame}
% =====================================================================

\subsection{Setting}
Let $D$ be a corpus of $N$ documents, $f_s$ the source (old) encoder, and $f_t$ the
target (new) encoder. An index $I_s$ over $\{f_s(d)\}_{d\in D}$ is serving. A query
workload $Q$ arrives over time, possibly with repeats. Let $\topk_K^{t}(q)$ denote the
native top-$K$ of $q$ under $f_t$ over all of $D$, and $\topk_{K'}^{s}(q)$ the
top-$K'$ window returned by $I_s$.

\subsection{Cost model}
The dominant cost of migration is encoding documents with $f_t$. We count
\emph{document encodings} $C$ as the primary cost unit. Query encoding, ANN search and
adapter application are orders of magnitude cheaper per call.

\textbf{Completion is a constraint, not a cost.} Every strategy we admit must end with
all $N$ documents encoded under $f_t$, so that the source index can be retired. This
has a consequence worth stating plainly, because it determines what the cost column can
mean: \emph{every admissible strategy encodes exactly $N$ documents}. We verified this
rather than assuming it --- of 828 completing strategy--order combinations in our runs,
none encodes fewer than $N$. Total encoding work therefore cannot separate strategies,
and any apparent saving in a reuse strategy is a matter of \emph{when} the work is
paid, not whether. What strategies do differ in is the quality they serve while the
migration runs, and that is what our objective measures.

\subsection{First-arrival measurement}
Quality is measured once per evaluation query, at the moment that query is first
served. A window strategy repairs its query's window inline and is charged for those
encodings at that moment; a blue--green rebuild serves the source index and is charged
for whatever background progress has been made. This avoids two errors that a
checkpoint-based replay invites: scoring unencoded window documents would let window
strategies reach target quality at $C{=}0$ without paying for the encodings, and
re-measuring the same query at successive checkpoints rewards strategies for being
measured late.

\textbf{Reference qualities.} $M_T$ is the quality (Recall@10 against qrels) of exact
native $f_t$ retrieval over all of $D$; $M_S$ is the same for $f_s$ --- what the system
serves before migration starts.

\textbf{Objective: debt.} For strategy $\sigma$, let $\bar M_\sigma$ be the mean quality
served over the evaluation population under first-arrival measurement. The objective is
\[
\debt(\sigma) = M_T - \bar M_\sigma ,
\]
the quality forgone during the migration. Lower is better; it can be negative, and is
on most of our pairs, because restricting the candidate set to a source window and
re-ranking it with the target model is a cascade that can exclude documents the target
model would itself have over-ranked.

\textbf{Admissibility.} A strategy--order combination is admissible on a pair if
(i) its background order completes, so the corpus reaches the target model;
(ii) it reaches target quality, $\bar M_\sigma \ge M_T - \varepsilon$; and (iii) it never
serves below the system it replaces, $\bar M_\sigma \ge M_S$. The floor is anchored to
$M_S$, not $M_T$: at $C{=}0$ every strategy serves roughly $M_S$, so a target-anchored
floor would exclude all strategies whenever the upgrade is worth doing.

\textbf{Winner rule, as pre-registered.} Among admissible combinations, those within an
equivalence band $\tau$ of the best $\debt$ are considered indistinguishable; within the
band we order by cost-at-service and then by a fixed simplicity ranking. We fixed
$\varepsilon = 0.01$, $\delta = 0.01$ and $\tau = 0.005$ on MS MARCO dev Recall@10
before any pair was run. Section~\ref{sec:select} reports the sensitivity of every
reported winner to $\tau$, because that parameter turns out to be load-bearing.

\textbf{Reporting cost.} We report two cost columns and they measure different things.
$C_{\text{total}}$ is the work a strategy ever does; as established above it is $N$ for
everything admissible. $C_{\text{serve}}$ is the work done by the time the strategy is
serving the evaluated population, and it is the quantity Section~\ref{sec:econ}
predicts. $C_{\text{serve}}$ is not a saving, and it is not one quantity: a strategy that
repaired windows reports work it did, whereas a strategy that reached target quality
having repaired nothing reports how far its background happened to get --- a number that
shrinks as traffic intensity rises while nothing about the strategy changes. Every
$C_{\text{serve}}$ we report carries a label saying which case produced it, and the two
are never compared with each other.

\subsection{Pair-level and workload-level measurements}
\emph{Compatibility} $\kappa = 1 - r$, where $r$ is the relative residual
$\lVert A(X_t) - X_s\rVert_F / \lVert X_s\rVert_F$ of a map $A$ from $f_t$ to $f_s$
embeddings (semi-orthogonal Procrustes), fit on one document sample and evaluated on a
disjoint one; larger $\kappa$ means more compatible. We also report linear
CKA~\cite{cka} on the same sample. \emph{Window gap}, following~\cite{embedflow},
$G(K') = M_T - M_{T\mid S_{K'}}$, where $M_{T\mid S_{K'}}$ is the quality of $f_t$
re-ranking the source window only. \emph{Residual miss}, its qrels-free counterpart,
\[
\rho(K') = 1 - \mathbb{E}_{q\in Q}\,
\frac{\lvert \topk_K^{t}(q) \cap \topk_{K'}^{s}(q)\rvert}{K}.
\]
Because $\rho$ depends on $K'/N$ we always report $K'$ as a fraction of $N$. Windows of
different depths are prefixes of a single retrieval at the largest depth, so $\rho(K')$
is non-increasing in $K'$ by construction; measurements use exact search and the
approximate-search noise floor is reported separately (pair P0).

\emph{Coverage}, the quantity Section~\ref{sec:econ} is built on:
\[
\cov(K') = \frac{1}{N}\Bigl\lvert \bigcup_{q \in Q_{\text{served}}} \topk_{K'}^{s}(q) \Bigr\rvert ,
\]
the fraction of the corpus that appears in at least one served window. It is computable
from the source index alone, with no target encoding at all.

% =====================================================================
\section{Strategies Under Study}\label{sec:strategies}
% =====================================================================

Table~\ref{tab:strategies} lists the strategies. Every strategy that does not reach the
target model on its own is paired with a background migration that re-encodes the corpus
at a fixed rate, so all admissible trajectories end at $M_T$ with $C = N$.

\begin{table*}[t]
\centering
\caption{Strategies evaluated. ``Nominator'' is what proposes candidates; ``Ranker''
is what orders them. S1 and S3 are the mapping family; S2, S4, S5 the window family.}
\label{tab:strategies}
\small
\begin{tabular}{@{}llll@{}}
\toprule
ID & Strategy & Nominator & Ranker \\
\midrule
S0 & Blue--green rebuild & --- (old index until switch) & $f_s$, then $f_t$ \\
S1 & Space mapping~\cite{drift-adapter} & $I_s$ via $A(f_t(q))$ & $f_s$ \\
S2 & Window re-rank~\cite{embedflow,thaw-anon} & $I_s$ via $f_s(q)$ & $f_t$ \\
S3 & Mapped window & $I_s$ via $A(f_t(q))$ & $f_t$ \\
S4 & Window $\cup$ lexical & $I_s \cup$ BM25~\cite{bm25} & $f_t$ \\
S5 & Window $\cup$ mapped $\cup$ lexical & all three & $f_t$ \\
\bottomrule
\end{tabular}
\end{table*}

\textbf{New writes.} Documents inserted during migration have no source vector and can
never be nominated by $I_s$. All window strategies therefore also search a small
target-space index over migrated and newly written documents and merge results.

\textbf{Background migration orders.} For S2--S5 the background re-encoder follows one
of: \emph{rand} (uniform random), \emph{size} (shortest first, maximizing documents per
GPU-second), \emph{demand} (no background work: only documents surfaced by queries are
encoded, which need not ever complete and is therefore inadmissible under
Section~\ref{sec:frame}), \emph{demand-then-rand} (demand first, random thereafter, which
does complete), and \emph{signal} (Section~\ref{sec:order}).

% =====================================================================
\section{Experimental Setup}\label{sec:setup}
% =====================================================================

\subsection{Corpora and queries}
\textbf{Text.} MS MARCO passages~\cite{msmarco}. The full-encoding experiments use a
$10^6$-passage subsample built as \emph{every judged passage} (MS MARCO dev and TREC-DL
2019/2020 qrels) plus a uniform random fill; a random subsample would drop most judged
passages and make $M_T$ and $G(K')$ meaningless. Recall@10 is computed on MS MARCO dev
(6{,}980 judged queries; primary).

\textbf{Code.} CodeSearchNet~\cite{codesearchnet} Python, 45{,}283 documents, with CoSQA
as an evaluation-only secondary query set.

\textbf{Workloads.} Uniform over the query set, and ORCAS~\cite{orcas} empirical
frequency for text. \todo{the reported runs use the uniform workload; add the ORCAS
and Zipf sweeps, or scope every workload-dependent claim to uniform}

\textbf{Query roles.} A seeded half of the judged queries (3{,}490 of 6{,}980) is held
out for evaluation; the rest form the traffic served during migration. Evaluated queries
enter the traffic only as explicit repeats. Without this split, on-demand repair warms
exactly the windows being evaluated and wins by construction; an earlier version of our
replay had this flaw and its results are discarded.

\textbf{Traffic intensity.} With background budget $b = N/R$ encodings per round, each
round serves $\lambda$ queries, set by the ratio $\mathrm{tr} = \lambda K'/b \in \{4, 16,
64, 256\}$: the on-demand encodings a fully cold round could trigger, relative to the
background budget. We report every setting and pre-register $\mathrm{tr}{=}64$ as the
primary. (We write this $\mathrm{tr}$, not $\tau$, which denotes the equivalence band.)

\subsection{Model pairs}
Table~\ref{tab:pairs} lists the pairs. P0 is the control: the same model with its
approximate index rebuilt under a different seed, so its $\rho(K')$ measures pure ANN
noise and serves as the noise floor. Pairs are ordered source\,$\to$\,target such that
$M_T > M_S$ on the held-out set, i.e.\ every pair is an upgrade someone would perform.

\begin{table*}[t]
\centering
\caption{Model pairs and per-pair measurements at $K' = 50$ ($K'/N = 5\times10^{-5}$ on
text). $\kappa$ is Procrustes compatibility, $\rho$ residual miss, $G$ the window gap,
$\cov$ union-of-windows coverage. $^\ddagger$P11's Procrustes residual exceeds 1, so the
orthogonal fit is worse than predicting zero and $\kappa$ is negative; it is a
compatibility \emph{failure}, not a small compatibility.}
\label{tab:pairs}
\small
\begin{tabular}{@{}llccccc@{}}
\toprule
ID & Source $\rightarrow$ Target & Type & $\kappa$ & CKA & $\rho(50)$ & $G(50)$ \\
\midrule
P0  & bge-base $\rightarrow$ same, ANN rebuilt & control      & $1.0000$ & $1.0000$ & $0.0000$ & $+0.0000$ \\
P2  & MiniLM-L6 $\rightarrow$ bge-base         & cross-family & $0.0337$ & $0.6112$ & $0.2064$ & $-0.0027$ \\
P3  & bge-small $\rightarrow$ bge-base         & same family  & $0.3876$ & $0.7668$ & $0.1395$ & $-0.0014$ \\
P4  & bge-m3 $\rightarrow$ bge-base            & same family  & $0.2875$ & $0.7191$ & $0.2234$ & $-0.0040$ \\
P5  & bge-m3 $\rightarrow$ bge-small           & same family  & $0.2737$ & $0.7103$ & $0.2412$ & $-0.0099$ \\
P9  & MiniLM-L6 $\rightarrow$ bge-small        & cross-family & $0.0419$ & $0.6510$ & $0.1871$ & $-0.0028$ \\
P10 & MiniLM-L6 $\rightarrow$ bge-m3           & cross-family & $0.0229$ & $0.5812$ & $0.2903$ & $-0.0094$ \\
P11 & MiniLM-L6 $\rightarrow$ e5-base          & cross-family & $-0.0928^\ddagger$ & $0.6032$ & $0.2060$ & $+0.0035$ \\
P12 & bge-m3 $\rightarrow$ e5-base             & cross-family & $0.3064$ & $0.7323$ & $0.2191$ & $-0.0023$ \\
P13 & bge-small $\rightarrow$ e5-base          & cross-family & $0.3985$ & $0.7474$ & $0.1684$ & $-0.0012$ \\
P14 & bge-base $\rightarrow$ e5-base           & cross-family & $0.4408$ & $0.8164$ & $0.1127$ & $-0.0007$ \\
\midrule
P7  & MiniLM-L6 $\rightarrow$ codesage-small   & general$\to$code & $0.0607$ & $0.3848$ & $0.3724$ & $+0.0110$ \\
\bottomrule
\end{tabular}
\end{table*}

\textbf{Pair coverage is BERT-lineage.} All eleven text pairs are encoder-only models in
the BERT lineage at 384, 768 or 1024 dimensions, and one pair crosses into a
code-specialized encoder. We have no cross-architecture pair (encoder\,$\to$\,LLM-based
embedder) and say so in Section~\ref{sec:disc}: it is the most consequential gap in this
study, because it is the only direction in which the compatibility axis might span a
usefully wider range. \todo{run the cross-family pair (RUNBOOK Phase 0.0) and either
widen this claim or record the prediction it confirms}

\subsection{Implementation and invariants}
Window, residual-miss and quality measurements use exact inner-product search, so they
are free of ANN error; HNSW~\cite{hnsw} with identical build parameters is used only for
P0 and for latency figures. Every encoding is computed once and cached, so all strategies
are replayed over identical vectors. Every model is encoded under an explicit contract
(prompt, pooling, padding side, EOS, revision) checked against the model card's worked
example and for batched-versus-single agreement; models that fail are excluded, not
tuned.

Before any pair is interpreted the replay must satisfy properties that hold by
construction: P0 with exact search gives $\rho(K'){=}0$; $\rho(K')$ is non-increasing in
$K'$; re-ranking a window equal to the corpus gives $G{=}0$; every admissible trajectory
ends at $M_T$ with $C{=}N$; S0 serves exactly $M_S$ until it switches; and the arrival
sequence is identical across strategies, so no strategy is compared on a different query
stream. A violation is a pipeline bug, not a finding. All twelve pairs satisfy all of
them.

% =====================================================================
\section{Pre-Registered Hypotheses and Their Outcomes}\label{sec:hyp}
% =====================================================================

We fixed five hypotheses, each with a falsification condition, before running any pair.
Table~\ref{tab:hyp} reports all five. Three are falsified on their stated criteria.
We report them as registered; the sections that follow give each one's measurement and
the mechanism behind the outcome.

\begin{table*}[t]
\centering
\caption{The five pre-registered hypotheses and their measured outcomes. Nothing in this
table was re-specified after seeing results.}
\label{tab:hyp}
\small
\begin{tabular}{@{}llp{0.40\textwidth}l@{}}
\toprule
 & Hypothesis & Falsification condition and measured value & Outcome \\
\midrule
H1 & $G(K')$ rises as $\kappa$ falls & Spearman $r_s < 0.5$ across pairs. Measured $r_s(\kappa, G) = +0.19$ over 11 non-control pairs --- below threshold and the wrong sign & \textbf{falsified} \\
H2 & $(\kappa,\rho)$ beats the majority baseline in LOPO mean regret & regret $\ge$ baseline's. Measured: majority $0.0018943$ at 9/12; $(\kappa,\rho)$ 1-NN $0.0020117$ at 8/12 & \textbf{falsified} \\
H3 & the low-$\kappa$/low-$\rho$ cell is non-empty (exploratory) & --- . P11 and P9 fall below the median on both axes; but $r_s(\kappa,\rho) = -0.52$, so the axes are far from independent & non-empty \\
H4 & the bounded-cost probe estimates $\rho(K')$ within $\pm 0.05$ at $\le 2\%$ of $N$ & error exceeds the bound on more than one pair. Measured: 7 of 11 non-control pairs exceed it at a $1.2\%$ total budget & \textbf{falsified} \\
H5 & the document signal has AUC $>0.5$ and \emph{signal} order beats \emph{rand} at equal budget & AUC CI excludes $0.5$ on 9 of 12 pairs ($0.625$--$0.652$). Head to head \emph{signal} wins 19 of 28 non-tied comparisons (80 of 108 tie): right in sign, sign-test $p = 0.087$, every effect inside $\tau$ & \textbf{inconclusive} \\
\bottomrule
\end{tabular}
\end{table*}

\note{F2, the result we lead with, was \emph{not} pre-registered. That coverage bounds
window cost from below holds by construction; the finding is that realised cost
\emph{attains} the bound, to four decimals, which we did not anticipate and report as a
post-hoc result. Section~\ref{sec:econ} is explicit about this.}

% =====================================================================
\section{Bounding the Cost of Windowed Reuse}\label{sec:econ}
% =====================================================================

This section is the result we would hand a practitioner, and the one we did not
anticipate.

\textbf{The bound.} To serve a query by re-ranking its source window, every document in
that window must be encoded under $f_t$. A window strategy serving query population
$Q_{\text{served}}$ must therefore encode at least $\bigl\lvert\bigcup_q
\topk_{K'}^{s}(q)\bigr\rvert = N\cdot\cov(K')$ documents before it is serving that
population at target quality, and a full rebuild costs exactly $N$. So no window strategy
can defer more than $1 - \cov(K')$ of the encoding work. This is arithmetic, not a
finding, and it is computable from the source index alone: one retrieval pass, no target
encoding, minutes rather than GPU-days.

\textbf{The finding: the bound is tight.} Table~\ref{tab:econ} compares the cost each
window strategy actually incurred by the time it was serving the evaluation population
against the statically computed $\cov(50)$ for its source encoder. The agreement is
$0.016$--$0.017\%$ relative, across eight pairs and three distinct source encoders.

\begin{table}[t]
\centering
\caption{The static bound is attained: realised cost-at-service against $\cov(50)$, at the
primary setting. $C_{\text{serve}}$ comes from the replay; $\cov(50)$ from a retrieval pass
that encodes nothing with $f_t$. A window strategy encodes the union of the windows it serves
by construction, so this table shows the bound is not exceeded by background work --- it is
not an independent prediction being confirmed. All eight rows have basis \emph{at-last-arrival}, so they are
mutually comparable (Section~\ref{sec:frame}).}
\label{tab:econ}
\small
\begin{tabular}{@{}llccc@{}}
\toprule
Source encoder & Pairs & $C_{\text{serve}}/N$ & $\cov(50)$ & rel.\ err. \\
\midrule
MiniLM-L6 & P2, P9, P10 & $0.264861$ & $0.264907$ & $0.017\%$ \\
bge-m3    & P4, P5, P12 & $0.255091$ & $0.255134$ & $0.017\%$ \\
bge-small & P3, P13     & $0.262135$ & $0.262176$ & $0.016\%$ \\
\bottomrule
\end{tabular}
\end{table}

\textbf{Why it is tight.} The realised cost converges on the \emph{whole-workload}
coverage, not the evaluation set's: by the time the last evaluation query is served, the
traffic queries have arrived and repaired their own windows, so the union of all served
windows has been encoded. The residual $0.017\%$ is the handful of documents in windows
whose queries had not yet arrived.

\textbf{What this buys.} The practitioner's question --- ``what will this migration cost
me, and when will I be serving the new model?'' --- is answerable before committing any
accelerator time, in closed form, from a quantity the source index already supports. No
replay, no target encoding, no simulation.

\textbf{Against the obvious baseline.} A reader should ask whether retrieval is needed at all:
if queries drew documents independently, coverage would be $1 - (1 - K'/N)^{|Q|}$, computable
from three integers. It is a real baseline and the measurement beats it, but not in the
direction one might guess --- the closed form \emph{overestimates} cost at every depth and every
encoder, because real windows overlap more than independent draws do:

\begin{center}
\small
\begin{tabular}{@{}rccr@{}}
\toprule
$K'$ & measured & closed form & over by \\
\midrule
50   & $0.2551$--$0.2663$ & $0.2946$ & $10.6$--$15.5\%$ \\
100  & $0.4138$--$0.4366$ & $0.5024$ & $15.1$--$21.4\%$ \\
200  & $0.6052$--$0.6453$ & $0.7524$ & $16.6$--$24.3\%$ \\
500  & $0.8346$--$0.8790$ & $0.9695$ & $10.3$--$16.2\%$ \\
1000 & $0.9392$--$0.9653$ & $0.9991$ & $3.5$--$6.4\%$ \\
\bottomrule
\end{tabular}
\end{center}

At $K' = 50$ the closed form predicts encoding $29.5\%$ of the corpus against a measured
$26.5\%$, understating the achievable saving by three points of $N$ --- 30{,}000 documents at
$N = 10^6$, and 107{,}000--147{,}000 at $K' = 200$, where the gap peaks. What the measurement
adds over the closed form is exactly the window concentration, and the practical reading is
that a planner using the closed form will under-provision the saving by $10$--$24\%$ in the
operating range. We report it because a cost model without its trivial baseline is not a
result.

\textbf{What the replay adds, and what it does not.} The agreement in
Table~\ref{tab:econ} is the weaker half of this section and we do not want it read as the
stronger: a window strategy encodes, by construction, the union of the windows it serves, so a
replay that reproduces that union to one document in 264{,}907 is confirming an
implementation. The substance is that the static quantity is a usable planning instrument and
that it beats the closed form; the replay's role is to show the bound is attained rather than
exceeded by background work.

\textbf{Coverage does not vary across encoders.} At $K' = 50$, coverage spans
$0.2551$--$0.2663$ over all five text encoders --- a range of $0.0112$. It is therefore
\emph{not} a usable third map axis (we had considered it as one); it is a property of the
corpus and workload, nearly independent of which encoder indexes them. That is good news
for the cost model and bad news for selection, and it foreshadows
Section~\ref{sec:select}.

\textbf{Window depth is a pure cost knob.} Coverage rises steeply with $K'$ --- $0.2649$,
$0.4374$, $0.6453$, $0.8787$, $0.9646$ at $K' = 50, 100, 200, 500, 1000$ --- while the
quality constraint never binds: the worst-case $G(K')$ over all eleven text pairs is
$+0.0035$ at $K' = 50$ and $\le +0.0000$ at every larger depth. Within the measured range
there is no quality price for a shallower window, so the depth should be set as small as
the constraint allows.

\todo{We measured $K' \ge 50$ only, so $K' = 50$ is the smallest depth we \emph{looked at}
and not the smallest \emph{admissible} one: the frontier is censored from below and the
$73.5\%$ ceiling reported here is a lower bound on what windowed reuse can defer. Run
$K' = 10, 25$ (RUNBOOK Phase 11.1) and either report a smaller admissible depth with a
larger ceiling, or report where $G$ first binds. Figure~\ref{fig:frontier} is blocked on
this.}

\todo{Figure~\ref{fig:frontier}: the coverage/$G$ frontier over the full $K'$ ladder, one
line per source encoder, with the admissibility threshold marked.}

% =====================================================================
\section{What the Strategies Cost}\label{sec:cost}
% =====================================================================

\textbf{Mapping never wins.} Sweeping the winner rule over every pair, every traffic
intensity and every $\tau \in [0, 0.01]$ --- 336 configurations --- the adapter-bridge
strategy S1 is selected $0$ times, the mapped window S3 $0$ times, and the union strategy
S5 $0$ times. Only S0 (39), S2 (276) and S4 (21) are ever selected. The ordering
$\text{S2} \approx \text{S4} \approx \text{S5} \gg \text{S0} \gg \text{S3} \gg \text{S1}$
holds on every pair.

The margin is not subtle. Bridged retrieval serves $M_A$ between $0.0514$ and $0.9497$
against target qualities of $0.83$--$0.88$; on the five worst pairs $M_A < 0.35$. In
$\debt$ terms S1 costs between $0.12$ and $0.82$ Recall@10, one to two orders of
magnitude worse than any window strategy, whose $\debt$ lies within $\pm 0.006$. A
practitioner choosing a bridge on the strength of its deferred encoding cost would be
trading almost all of their retrieval quality for it.

\textbf{$\kappa$ predicts how badly mapping fails.} The compatibility axis does not
predict the winner (Section~\ref{sec:select}), but it does predict the adapter's served
quality: $r(\kappa, M_A) = +0.742$ over the ten text pairs excluding the control
($n = 10$, $p \approx 0.014$). Linear CKA ties it on that subset ($+0.739$) but its
apparent signal over all twelve pairs is carried entirely by the identity control ---
dropping P0 alone takes $r(\mathrm{CKA}, M_A)$ to $-0.122$ while $\kappa$ holds at
$+0.415$. $\rho$ has the wrong sign for a compatibility measure, $-0.572$. We report this
as a secondary observation at $n = 10$, not an established predictor, and note it does
not transfer across corpora: the code pair has $\kappa = 0.0607$ and the highest $M_A$ of
any pair, $0.9497$, because its source model is already strong on that corpus.

\textbf{``Do nothing'' is a tautology, not a strategy.} On eleven of twelve pairs the
blue--green rebuild's $\debt$ equals the quality gap $M_T - M_S$ \emph{exactly} --- to
$0.00\mathrm{e}{+}00$ --- because it serves source quality throughout the measured period.
So S0 is admissible precisely when the upgrade is worth less than the equivalence band,
and the two pairs on which it is selected (P0 at $0.0000$ and P14 at $0.0043$, against
$\tau = 0.005$) are exactly the two whose gap falls inside it. This is worth stating
because it looks like a strategy result and is not: it is a restatement of the quality
gap. Every other pair has a gap of at least $0.0100$.

\todo{Table: $\debt$ by strategy family per pair, from results/tables/grid.csv, with
tied-set sizes.}

% =====================================================================
\section{Selection is Moot}\label{sec:select}
% =====================================================================

H2 asked whether $(\kappa, \rho)$ predicts the winner on a held-out pair. We hold out
each pair in turn, fit a decision rule on the remaining pairs, and predict the held-out
winner, scoring by regret in the objective's units.

\begin{table}[t]
\centering
\caption{Leave-one-pair-out winner prediction. The majority rule uses no features at all.
Mean regret is in Recall@10. The two feature rules that match the baseline's accuracy
reproduce its regret to every digit, i.e.\ they are the constant rule in disguise.}
\label{tab:lopo}
\small
\begin{tabular}{@{}lcc@{}}
\toprule
Rule & Accuracy & Mean regret \\
\midrule
\textbf{majority (no features)} & \textbf{9/12} & \textbf{0.0018943} \\
1-NN on $\kappa$ only          & 9/12 & 0.0018943 \\
depth-2 tree on $\kappa$ only  & 9/12 & 0.0018943 \\
1-NN on $(\kappa, \rho)$       & 8/12 & 0.0020117 \\
depth-2 tree on $(\kappa, \rho)$ & 8/12 & 0.0028713 \\
depth-2 tree on $\rho$ only    & 8/12 & 0.0028713 \\
1-NN on $\rho$ only            & 6/12 & 0.0032533 \\
\bottomrule
\end{tabular}
\end{table}

H2 is falsified: the $(\kappa,\rho)$ rule's mean regret, $0.0020117$, exceeds the
majority baseline's $0.0018943$. Every rule that actually discriminates on features does
worse than ignoring them, and $\rho$ alone is near chance.

\textbf{Why, mechanically.} Three reasons, each measurable.
(i) \emph{There is almost nothing to predict.} At the primary setting the twelve labels
are S2 nine times, S0 twice --- both tautologies, as Section~\ref{sec:cost} showed --- and
S4 once. No feature can separate what does not vary.
(ii) \emph{At the pre-registered traffic intensity the window variants are
indistinguishable.} This is a property of that operating point, not of the method, and we
were initially wrong to state it as one: window strategies also search the target-side index
over already-migrated documents, so what the background has reached does affect the served
answer. It has simply reached almost nothing by the time a dense workload is served ---
cost-at-service for the rebuild is $0.03$ of $N$ at $\mathrm{tr}{=}64$ against $0.63$ at
$\mathrm{tr}{=}4$ --- and at $\mathrm{tr} \in \{4, 16\}$ the order does move $\debt$, on
every non-degenerate pair. Eight pairs reported
one of 3--4 combinations with identical $\debt$, and the control one of fourteen spanning
all six families; a reported winner such as ``S2/demand-then-rand'' is one arbitrary
member of a tie, chosen by our simplicity rule.
(iii) \emph{The baseline is already at the achievable ceiling.} The constant rule's mean
regret, $0.0019$, is \emph{below} the equivalence band $\tau = 0.005$ that we declared
before running anything. Its mistakes fall inside the region the study itself calls
indistinguishable.

\textbf{Sensitivity to $\tau$.} The two pairs that make the map look non-trivial are
exactly the two that move with the band. At $\tau \le 0.004$, P14 selects the window
family; at $\tau \ge 0.006$, P11 does; $5$ of $48$ pair\,$\times$\,setting cells differ
from the registered $\tau$ in one direction and $4$ at the other. At $\tau = 0$ the band
never engages and $22$ of $48$ differ, with selection falling to $\debt$ differences of
order $10^{-4}$. So one or two pairs always deviate from S2, but \emph{which} pair is a
function of $\tau$ rather than of the pair. We publish the full curve rather than a single
column, and we did not retune $\tau$ after seeing it.

\todo{Figure: winner by pair against $\tau$, from results/tables/tau\_curve.csv.}

% =====================================================================
\section{A Bounded-Cost Probe}\label{sec:probe}
% =====================================================================

Measuring $\rho(K')$ directly requires the native target top-$K$, i.e.\ encoding all of
$D$ with $f_t$ --- the very cost migration seeks to avoid. H4 asked whether it can be
estimated at bounded cost.

\textbf{Intruder-count estimator.} For a probe query $q$, encode the window
$\win_q = \topk_{K'}^{s}(q)$ with $f_t$ and let $w_1 \ge \dots \ge w_K$ be the top-$K$
target scores in $\win_q$. A document outside the window can displace a window document
from the target top-$K$ only if its target score exceeds some $w_j$. With
$I_q(w) = \lvert\{d\notin\win_q : s_t(q,d) > w\}\rvert$, the overall top-$K$ contains $j$
outside documents iff the $j$-th best outside score exceeds $w_{K-j+1}$, so
\[
\mathrm{miss}_q = \bigl\lvert\{\, j \in 1..K : I_q(w_{K-j+1}) \ge j \,\}\bigr\rvert .
\]
Drawing a uniform sample $S\subset D$ of size $m$ and encoding it with $f_t$ gives the
unbiased count estimate $\hat I_q(w) = (N/m)\lvert\{d \in S\setminus \win_q : s_t(q,d) >
w\}\rvert$, which we plug in. One sample serves every probe query, and every encoding it
performs counts toward the migration, so none is wasted.

\textbf{Result: H4 is falsified, but the stratified estimator is the right design.}
Intruders are rare tail events, so the uniform estimator has high variance. Stratifying
on source rank --- oversampling documents at source ranks $K' < r \le cK'$, with
importance weights --- cuts the mean absolute error from $0.1524$ to $0.0573$ at a $1\%$
sample, a factor of $2.7$. But the pre-registered bound was $\pm 0.05$ with falsification
on more than one pair, and $7$ of $11$ non-control pairs exceed it at a $1.2\%$ total
budget. H4 fails as stated.

\textbf{Diagnosis.} The binding constraint is the number of probe queries, not the sample
size. The window term costs $\lvert Q_{\text{probe}}\rvert \cdot K'$ encodings, so the
budget admitted only 12 probe queries; $\rho$ estimated from 12 queries has sampling error
of roughly the size of the miss itself. This is a budget allocation we would set
differently --- more queries, shallower windows --- but we report the registered design.
Since $\rho$ turns out not to predict the winner (Section~\ref{sec:select}), a better
probe for it would not change any recommendation in this paper, and we do not pursue it.

\begin{table}[t]
\centering
\caption{Probe error at a $1\%$ sample ($\approx 1.2\%$ total budget, 12 probe queries,
$K' = 50$). H4's bound is $\pm 0.05$; rows above it are marked.}
\label{tab:probe}
\small
\begin{tabular}{@{}lcccc@{}}
\toprule
Pair & $\rho$ true & uniform err. & stratified err. & vs.\ bound \\
\midrule
P0  & $0.0000$ & $0.0000$ & $0.0000$ & pass \\
P3  & $0.1417$ & $0.1333$ & $0.0167$ & pass \\
P4  & $0.1833$ & $0.1083$ & $0.0250$ & pass \\
P12 & $0.2083$ & $0.1250$ & $0.0500$ & pass \\
P14 & $0.1250$ & $0.1167$ & $0.0500$ & pass \\
P2  & $0.2083$ & $0.2083$ & $0.0667$ & exceeds \\
P5  & $0.2000$ & $0.1500$ & $0.0667$ & exceeds \\
P10 & $0.2833$ & $0.2250$ & $0.0667$ & exceeds \\
P9  & $0.1917$ & $0.1917$ & $0.0833$ & exceeds \\
P13 & $0.2333$ & $0.2167$ & $0.0833$ & exceeds \\
P7  & $0.1375$ & $0.1375$ & $0.0875$ & exceeds \\
P11 & $0.2667$ & $0.2167$ & $0.0917$ & exceeds \\
\bottomrule
\end{tabular}
\end{table}

% =====================================================================
\section{Document-Level Migration Ordering}\label{sec:order}
% =====================================================================

\textbf{Signal.} Some documents are retrieved under the source model because of lexical
evidence; others only because of the source model's semantics. We hypothesized the latter
are the ones whose exposure changes under the target model. With
$\mathrm{kocc}_s(d)$ the number of workload queries for which $d\in\topk_K^{s}(q)$, the
\emph{dense-only heat} is
\[
h(d) = \bigl\lvert\{q : d\in\topk_K^{s}(q) \wedge d\notin\topk_{K_b}^{\mathrm{BM25}}(q)\}\bigr\rvert ,
\]
computable from the old index and a BM25 index before any target encoding. The gate uses
the dense-only share $h(d)/\mathrm{kocc}_s(d)$ against the label ``relative exposure
change exceeds its median''.

\textbf{The signal is real.} The bootstrap $95\%$ CI of the AUC excludes $0.5$ on nine of
twelve pairs, with AUC between $0.6254$ and $0.6515$ on those nine. The two pairs where it
fails are informative: P10 ($0.4546$, CI $[0.4181, 0.4912]$) and the code pair P7
($0.4676$, CI $[0.4625, 0.4729]$) are both below chance.

\textbf{And the head-to-head is inconclusive.} H5 also required \emph{signal} order to beat
\emph{rand} at equal budget. Across all pairs, traffic intensities and window strategies,
\emph{signal} is strictly better in 19 comparisons and strictly worse in 9, with 80 of 108
tied: right in sign, but a sign test on the 28 non-tied comparisons gives $p = 0.087$, and the
largest single effect is $0.0013$ Recall@10 against an equivalence band of $0.005$. The
ordering helps, consistently in direction, by an amount this study has declared
indistinguishable. We report H5 as inconclusive rather than positive or negative.

The ties fall where one would expect. At $\mathrm{tr} \in \{64, 256\}$ the evaluation
population is served before the background has migrated materially, so the order cannot
matter, and all 80 ties are there. At $\mathrm{tr} \in \{4, 16\}$, where the background has
covered a large share of the corpus by service time, the order does move $\debt$ --- and every
non-tied comparison lives there. An ordering signal pays only in proportion to how much
migration happens while the system is still being measured, which under a dense workload and
a completion requirement is a narrow window. That also bounds how much RQ4 could ever have
been worth in this setting, and we would say so in the design rather than in the results if we
were starting again.

% =====================================================================
\section{The Regime Where None of This Helps}\label{sec:regime}
% =====================================================================

Windowed reuse defers $1 - \cov(K')$ of the encoding work, and
$\cov(K') \to 1$ as $\lvert Q\rvert K'$ approaches $N$. Our code corpus sits on the far
side of that boundary: 43{,}730 queries over 45{,}283 documents, giving
$\cov(K') = 1.0000$ at every depth from 23 to 1000. The union of the windows is the
corpus, the ceiling on deferral is exactly zero, and no window strategy can save anything
there however compatible the models are.

This is a property of the corpus and workload, not of the models, and it is checkable in
advance by the same static measurement as Section~\ref{sec:econ}. It is also why we do
\emph{not} present the code pair as cross-domain evidence about strategy selection: where
only one strategy is ever economic, the selection question is vacuous. P7's value is as a
boundary case: the bound degenerates correctly at $\cov = 1$, where it says reuse can buy
nothing.

\todo{One code pair is a thin basis for a regime claim. Either add a code$\to$code pair
(codesage-small $\to$ codesage-base) or scope this section explicitly to one corpus.}

% =====================================================================
\section{Discussion and Limitations}\label{sec:disc}
% =====================================================================

\textbf{What a practitioner should do.} Measure $\cov(K')$ on the source index at the
depth you would use --- one retrieval pass, no target encoding. If
$1 - \cov(K')$ is small, windowed reuse cannot pay and a blue--green rebuild is the
honest choice. If it is large, use plain window re-ranking (S2) with any completing
background order; do not spend effort choosing among the window variants, and do not use
an adapter bridge. If $M_T - M_S$ is smaller than the quality regression you are willing
to absorb, do not migrate at all.

\textbf{What we got wrong.} We designed this study around a predictor and it does not
exist in the range we measured. We think the honest reading is not ``compatibility is
irrelevant'' but ``the strategy space collapses before the compatibility axis gets a
chance to matter'': one family dominates, its variants are provably tied, and the only
genuine fork --- migrate or do not --- is decided by a quantity ($M_T - M_S$) that needs
no map. A map needs labels that vary, and ours do not.

\textbf{Threats to validity.}
\begin{itemize}
\item \emph{Pair coverage.} All eleven text pairs are BERT-lineage encoder-only models.
We have no encoder\,$\to$\,LLM-embedder pair, which is the direction most likely to widen
the compatibility range. Our results are evidence about these pair types, not a universal
law, and we state a falsifiable prediction for that case: since a bridge needs
$M_A \ge M_T - \tau \approx 0.870$ and the best we measured is $0.757$ at $\kappa =
0.4408$, and cross-architecture $\kappa$ should run \emph{lower} than within-lineage
$\kappa$, we predict mapping fails there too. \todo{run it}
\item \emph{One code corpus, one code pair.}
\item \emph{One workload.} The reported runs use the uniform workload; the ORCAS and Zipf
sweeps are specified but not run, so every claim here is scoped to uniform arrival.
\todo{run them or scope the claims in each section}
\item \emph{Subsampled corpora.} $\rho$ and $\cov$ depend on $K'/N$; we report $K'$ as a
fraction and note that $10^6$ is two to three orders of magnitude below production scale.
\item \emph{The window depth is not justified from below.} See the \todo{} in
Section~\ref{sec:econ}: $K' = 50$ is the smallest depth measured, not the smallest
admissible one.
\item \emph{$\tau$ is a judgement.} We fixed $0.005$ before running and publish the
sensitivity curve rather than claiming the value is principled.
\item \emph{Encoding cost as the cost unit} understates query-time latency for window
strategies, which encode documents inline on a cold query; we report encodings per served
query separately.
\item \emph{Scope.} We study complete migration only. The strategy that dominates on both
axes in our own data --- lazy materialization, which never completes --- is excluded by
our admissibility constraint by construction, and its economics are concurrent
work~\cite{thaw-anon}. A reader should not conclude from this paper that completion is
the right choice, only that under completion the reuse question has the shape described
here.
\end{itemize}

% =====================================================================
\section{Related Work}\label{sec:related}
% =====================================================================

\textbf{Embedding model upgrades.} Drift-Adapter maps new queries into the legacy space
to defer re-encoding~\cite{drift-adapter}; our S1/S3 are its families, and
Section~\ref{sec:cost} is in effect a negative replication at $10^6$ scale under a
quality-along-the-path objective. EmbedFlow serves a target model over source-index
candidates while materializing target vectors, with an empirical window diagnostic and a
registry of measured pairs~\cite{embedflow}; our $G(K')$ follows its definition and our S2
is its mechanism. Residual reachability reuses old graph structure to accelerate
new-index construction~\cite{dgm}. Concurrent work studies exactness conditions and the
economics of incomplete migration~\cite{thaw-anon}; we deliberately measure the
complementary complete-migration regime. Our contribution is the common frame, the static
cost model, and the negative selection result.

\textbf{Representation alignment.} Linear and orthogonal alignment of embedding
spaces~\cite{procrustes}, representation similarity~\cite{cka}, and relative
representations~\cite{relrep} motivate our compatibility axis. Our finding that linear
CKA's apparent predictive signal is carried by an identity control is a caution for its
use as a practical compatibility score.

\textbf{Dynamic vector indexes.} SPFresh, FreshDiskANN, and Quake update indexes as
vectors change~\cite{spfresh,freshdiskann,quake}; they assume the new vector exists, which
is the cost studied here. LEANN recomputes embeddings at query time to save
storage~\cite{leann}, a mechanism related to window re-ranking with a different goal.

\textbf{Stale representations.} Corrector networks reduce re-encoding of stale embeddings
during retriever training~\cite{corrector}; CIPBench studies which content changes require
re-embedding under a fixed encoder~\cite{cipbench}, and reports that an exposure signal
added nothing over change history --- consistent with our Section~\ref{sec:order}.

% =====================================================================
\section{Conclusion}
% =====================================================================

We evaluated three families of embedding-upgrade reuse strategies over twelve model pairs
under one cost model. We set out to build a predictor and found there is little to
predict: one family dominates, its variants are provably tied, and three of our five
pre-registered hypotheses failed on their own criteria. What the measurements do deliver
is an economic account that is more directly useful than the map we wanted. The cost of
windowed reuse is bounded below by the union of the windows a workload touches; that union is
computable from the source index in minutes, is $10$--$24\%$ tighter than the closed form that
assumes independent windows, and the replay confirms it is attained. The upgrade decision is
therefore a regime question, answerable in advance from a bound, rather than a
strategy-selection question. Mapping, meanwhile, does not work at
this scale: it never wins, and the compatibility score tells you in advance how badly it
will lose.

\section*{Artifacts}
\todo{Repository URL (anonymized): replay framework, configs, encoder contracts, golden
files, coverage and probe outputs, all per-pair result JSONs and the derived tables. MS
MARCO- and ORCAS-derived vectors are not redistributable; scripts regenerate them.}

% =====================================================================
\begin{thebibliography}{99}
\bibitem{drift-adapter} H.~Vejendla, ``Drift-Adapter: A practical approach to near
zero-downtime embedding model upgrades in vector databases,'' in \emph{Proc. EMNLP},
2025, pp.~15938--15949.
\bibitem{embedflow} A.~Srivastav, ``EmbedFlow: Progressive embedding-model migration
over existing vector indexes,'' software, v0.8.0, 2026.
\url{https://github.com/arnsri33/embedflow}
\bibitem{dgm} \todo{authors}, ``Exploiting residual reachability for cross-model
migration of graph-based indexes in approximate nearest neighbor search,''
arXiv:2609.23036, 2026.
\bibitem{thaw-anon} Anonymous, ``\todo{Thaw title},'' under submission, 2026.
\bibitem{cipbench} N.~Zhao, J.~Zhang, T.~Zhong, H.~Wang, and B.~Wang, ``Which code
changes need re-embedding? CIPBench and the limits of multi-signal index
maintenance,'' \emph{ACM TOSEM} (submitted), 2026.
\bibitem{zilliz-migration} Zilliz, ``Migrate to a new embedding model,'' documentation,
2026.
\bibitem{msmarco} T.~Nguyen \emph{et al.}, ``MS MARCO: A human generated machine reading
comprehension dataset,'' arXiv:1611.09268, 2016.
\bibitem{orcas} N.~Craswell \emph{et al.}, ``ORCAS: 18 million clicked query-document
pairs for modeling search,'' in \emph{Proc. CIKM}, 2020.
\bibitem{codesearchnet} H.~Husain \emph{et al.}, ``CodeSearchNet challenge: Evaluating
the state of semantic code search,'' arXiv:1909.09436, 2019.
\bibitem{bm25} S.~Robertson and H.~Zaragoza, ``The probabilistic relevance framework:
BM25 and beyond,'' \emph{Found. Trends Inf. Retr.}, vol.~3, no.~4, 2009.
\bibitem{hnsw} Y.~Malkov and D.~Yashunin, ``Efficient and robust approximate nearest
neighbor search using hierarchical navigable small world graphs,'' \emph{IEEE TPAMI},
vol.~42, no.~4, 2018.
\bibitem{cka} S.~Kornblith, M.~Norouzi, H.~Lee, and G.~Hinton, ``Similarity of neural
network representations revisited,'' in \emph{Proc. ICML}, 2019.
\bibitem{procrustes} P.~Sch\"onemann, ``A generalized solution of the orthogonal
Procrustes problem,'' \emph{Psychometrika}, vol.~31, 1966.
\bibitem{relrep} L.~Moschella \emph{et al.}, ``Relative representations enable zero-shot
latent space communication,'' in \emph{Proc. ICLR}, 2023.
\bibitem{hubness} M.~Radovanovi\'c, A.~Nanopoulos, and M.~Ivanovi\'c, ``Hubs in space:
Popular nearest neighbors in high-dimensional data,'' \emph{JMLR}, vol.~11, 2010.
\bibitem{spfresh} Y.~Xu \emph{et al.}, ``SPFresh: Incremental in-place update for
billion-scale vector search,'' in \emph{Proc. SOSP}, 2023.
\bibitem{freshdiskann} A.~Singh \emph{et al.}, ``FreshDiskANN: A fast and accurate
graph-based ANN index for streaming similarity search,'' arXiv, 2021.
\bibitem{quake} J.~Mohoney \emph{et al.}, ``Quake: Adaptive indexing for vector search,''
in \emph{Proc. OSDI}, 2025.
\bibitem{leann} \todo{authors}, ``LEANN: A low-storage vector index,'' arXiv, 2025.
\bibitem{corrector} N.~Monath \emph{et al.}, ``A fresh take on stale embeddings:
Improving dense retriever training with corrector networks,'' in \emph{Proc. ICML},
2024.
\end{thebibliography}

\end{document}
